%%%%%%%%%%%%%%%%%%%%%%%%%%%%%%%%%%%%%%%%
\chapter{Resumen}
%%%%%%%%%%%%%%%%%%%%%%%%%%%%%%%%%%%%%%%%
Esta monografía presenta una revisión sistemática de las técnicas de \textit{fine-tuning} para \textit{Large Language Models} (LLMs) publicadas entre 2020 y 2025, conducida mediante el protocolo PRISMA 2020. A partir de una búsqueda inicial de 140 registros en Scopus, se seleccionaron 84 estudios empíricos con evaluación cuantitativa, distribuidos en cuatro dominios: Salud y Biomedicina (35,7\%), Lenguaje y Humanidades (25,0\%), Ciencias Exactas y Computación (23,8\%) e Industria y Negocios (15,5\%). El análisis abarcó nueve técnicas organizadas en tres paradigmas: \textit{Full Fine-Tuning} y SFT como enfoques de actualización completa de parámetros, cinco métodos de \textit{Parameter-Efficient Fine-Tuning} (LoRA, QLoRA, \textit{Adapter Layers}, \textit{Prefix Tuning}, \textit{Prompt Tuning}) y dos métodos de alineación (RLHF, DPO). Los resultados evidencian la hegemonía de LoRA y sus variantes, presente en el 79,8\% de los estudios incluidos, con una calidad metodológica del corpus aceptable (\textit{score} promedio de 2,61 sobre 3,0). Se identificaron tres patrones transversales: la efectividad consistente del \textit{fine-tuning} sobre el modelo base sin adaptar, la ventaja de las técnicas PEFT en lenguas y dominios de baja cobertura en el preentrenamiento, y la tendencia hacia arquitecturas compuestas en aplicaciones de producción. Los hallazgos permiten formular criterios de selección basados en evidencia empírica según el contexto operativo, las restricciones de privacidad y los requisitos de alineación.

\paragraph{Palabras clave:} \textit{fine-tuning}, \textit{Large Language Models}, LoRA, \textit{Parameter-Efficient Fine-Tuning}, RLHF, DPO, revisión sistemática, PRISMA, procesamiento de lenguaje natural, evaluación empírica.

%%%%%%%%%%%%%%%%%%%%%%%%%%%%%%%%%%%%%%%%
\chapter{Abstract}
%%%%%%%%%%%%%%%%%%%%%%%%%%%%%%%%%%%%%%%%
This monograph presents a systematic review of fine-tuning techniques for Large Language Models (LLMs) published between 2020 and 2025, conducted following the PRISMA 2020 protocol. From an initial search of 140 records in Scopus, 84 empirical studies with quantitative evaluation were selected, distributed across four application domains: Health and Biomedicine (35.7\%), Language and Humanities (25.0\%), Exact Sciences and Computing (23.8\%), and Industry and Business (15.5\%). The analysis covered nine techniques organized into three paradigms: Full Fine-Tuning and SFT as full-parameter update approaches, five Parameter-Efficient Fine-Tuning methods (LoRA, QLoRA, Adapter Layers, Prefix Tuning, Prompt Tuning), and two alignment methods (RLHF, DPO). Results reveal the dominance of LoRA and its variants, present in 79.8\% of included studies, with an acceptable methodological quality of the corpus (average score of 2.61 out of 3.0). Three cross-domain patterns were identified: the consistent effectiveness of fine-tuning over the unadapted base model, the advantage of PEFT techniques in low-resource languages and underrepresented domains, and the trend toward composite architectures in production applications. The findings enable evidence-based technique selection criteria according to operational context, data privacy constraints, and response alignment requirements.

\paragraph{Keywords:} fine-tuning, Large Language Models, LoRA, Parameter-Efficient Fine-Tuning, RLHF, DPO, systematic review, PRISMA, natural language processing, empirical evaluation.